\documentclass[10pt]{article}
\usepackage[margin=0.6in,top=0.55in,bottom=0.55in]{geometry}
\usepackage{graphicx}
\usepackage{amsmath,amssymb}
\usepackage[hidelinks]{hyperref}
\usepackage{titlesec}
\usepackage[font=small,labelfont=bf]{caption}
\usepackage{mathptmx}  % Times New Roman for text and math

% Compact sections
\titleformat{\section}{\normalsize\bfseries}{\thesection}{0.5em}{}
\titlespacing*{\section}{0pt}{7pt}{3pt}
\setlength{\parskip}{3pt}
\setlength{\parindent}{0pt}
\setlength{\abovecaptionskip}{4pt}
\setlength{\belowcaptionskip}{2pt}
\setlength{\textfloatsep}{6pt}
\setlength{\floatsep}{6pt}
\setlength{\intextsep}{6pt}

\graphicspath{{../results/manuscript_figures/}{../results/figures/}}

\begin{document}

\begin{center}
{\Large\bfseries DNA foundation models encode motif identity\\[2pt] but fail to learn regulatory grammar}
\vspace{5pt}
\end{center}

\section*{Background}

Whether enhancer activity is governed by motif arrangement or by more permissive logic remains a central question in cis-regulatory biology~[1,2].
We propose the \textbf{billboard model}: DNA foundation models encode which transcription factor binding sites are present (motif identity) but not their relative positions, orientations, or spacing (motif arrangement).
Under this model, enhancers function like billboards---what matters is which elements appear, not how they are arranged.
This is distinct from asking whether biology itself follows billboard logic; we specifically test what models learn.
We derive testable predictions and evaluate them using a spacer-factored grammar sensitivity framework (SF-GSI) that isolates motif arrangement effects from motif identity effects---resolving a critical confound in standard vocabulary-preserving shuffles, which perturb both simultaneously.
Crucially, we also test whether grammar matters biologically by identifying MPRA sequence pairs with identical motif vocabulary but different arrangements.

\begin{figure}[h!]
\centering
\includegraphics[width=\textwidth,keepaspectratio]{fig1_spacer_confound.pdf}
\caption{\textbf{Absolute context dominates apparent grammar effects.}
\textbf{(A)}~Factorial decomposition of prediction variance into position, orientation, and spacer components across three MPRA datasets (DNABERT-2); spacer composition consistently accounts for the largest share.
\textbf{(B)}~Spacer ablation series: replacing spacer DNA produces larger expression changes than motif-only rearrangement.
\textbf{(C)}~Feature-based regression predicting model expression output from compositional sequence features; compositional features alone explain 40--80\% of variance.
\textbf{(D)}~GC--expression correlation reverses sign across datasets.}
\label{fig:spacer}
\end{figure}

\section*{Methods}

SF-GSI computes a grammar sensitivity index, $GSI = \sigma(\hat{y}) / |\mu(\hat{y})|$, over vocabulary-preserving motif shuffles and factors the response into positional versus absolute context using matched-spacer perturbations and null-normalized testing.

The framework proceeds in three stages: (1) factorial decomposition isolating contributions of motif position, orientation, and spacer composition; (2) spacer ablation series applying graded compositional perturbations; (3) feature decomposition regressing model predictions against explicit compositional features.
Statistical significance is assessed via a three-stage correction cascade: F-test, z-score null normalization, and Benjamini--Hochberg FDR correction at $\alpha = 0.05$.

We applied SF-GSI to five MPRA datasets (Agarwal, Klein, Inoue, Vaishnav, and Jores), totaling 7,650 model--enhancer measurements across DNABERT-2, Nucleotide Transformer v2, HyenaDNA, Enformer, and PARM~[3--8]. Code: \url{https://github.com/bryanc5864/GRAMLANG}.

\begin{figure}[h!]
\centering
\includegraphics[width=\textwidth,keepaspectratio]{fig3_gsi_census.pdf}
\caption{\textbf{Statistical correction collapses apparent grammar prevalence.}
\textbf{(A)}~Violin plots of raw GSI values by model.
\textbf{(B)}~Median GSI profiles across five datasets.
\textbf{(C)}~Correction cascade: grammar-sensitive calls fall from 100\% to 8.3\% after null normalization and to 0.12\% (9/7,650) after FDR correction.
\textbf{(D)}~Pairwise cross-model GSI correlations by dataset.}
\label{fig:census}
\end{figure}

\section*{Results}

\textbf{Models fit the billboard model.} Under naive thresholding all tested enhancers appear grammar-sensitive, but this falls to 8.3\% after null normalization and to 0.12\% (9/7,650) after FDR correction.
Per-enhancer classification shows 89.7\% (95\% CI: 87.7--91.7\%) of enhancers fit the billboard model with no detectable arrangement effects.
Spacer composition accounts for 78--86\% of apparent grammar-sensitivity variance, while motif-only rearrangements produce 2--6$\times$ smaller effects than spacer perturbations.

\textbf{But grammar matters biologically.} We identified 363 MPRA sequence pairs with identical motif vocabulary but different arrangements.
These pairs show measured expression differences up to $\Delta = 6.3$ (Jores dataset), with 76\% having $|\Delta| > 0.5$.
This is direct experimental evidence that arrangement affects expression---grammar exists.

\textbf{Models fail to learn grammar.} We tested whether models can detect grammar when it demonstrably exists: (1) synthetic sequences with obvious grammar rules yielded $p = 0.73$---not detected; (2) known helical periodicity (10.5bp) yielded $p = 0.96$---not detected.
Models cannot learn even simple, known grammar patterns.

Positive controls confirm genuine grammar is detectable with controlled spacers ($p \le 7.2\times10^{-15}$), but this requires experimental designs that current models were not trained on.

\section*{Conclusions}

These results support the billboard model: DNA foundation models encode motif identity but fail to learn regulatory grammar.
We found 363 sequence pairs in MPRA data with identical motif vocabulary but different arrangements that show measured expression differences up to $\Delta = 6.3$, proving grammar matters biologically.
Yet models cannot detect even synthetic grammar rules ($p = 0.73$) or known helical periodicity ($p = 0.96$).
The billboard model (89.7\% of enhancers, 95\% CI: 87.7--91.7\%) describes model behavior, not biological reality---current architectures systematically miss regulatory grammar that demonstrably exists.
This reframes model evaluation: benchmarks should test whether models can learn grammar that exists in experimental data.

\vspace{3pt}
\noindent\textbf{Availability:} \url{https://github.com/bryanc5864/GRAMLANG}

\vspace{4pt}
\noindent\textbf{\small References}\\
{\footnotesize
\noindent [1]~Arnosti \& Kulkarni, \textit{J.~Cell.~Biochem.} 94:890, 2005.
\enspace [2]~Long et~al., \textit{Cell} 167:1170, 2016.\\
\noindent [3]~Zhou et~al., \textit{ICLR} 2024, DNABERT-2.
\enspace [4]~Dalla-Torre et~al., \textit{Nature Methods} 2024, Nucleotide Transformer.\\
\noindent [5]~Nguyen et~al., \textit{arXiv}:2306.15794, 2023, HyenaDNA.
\enspace [6]~Avsec et~al., \textit{Nature Methods} 18:1196, 2021, Enformer.\\
\noindent [7]~de~Jonge et~al., \textit{Nature} 2025, PARM.\\
\noindent [8]~Agarwal et~al., \textit{Nature} 639:411, 2025; Jores et~al., \textit{Nat.~Plants} 7:842, 2021.
}

\end{document}
